\begin{abstract}
\noindent Le \emph{network slicing} partage une infrastructure 5G/B5G entre des services aux exigences opposées (eMBB, URLLC, mMTC). L'opérateur doit réserver à chaque slice une part des ressources suffisante pour son contrat, et vérifier cette allocation quand le réseau change brusquement. À partir du dataset \emph{B5G Network Slicing} (Farreras et al., 2024 : 7\,127 simulations paquet, 1,58 million de flux, 162 topologies réelles), nous formulons le problème comme un \textbf{jumeau numérique}. Il s'agit de prédire, pour chaque flux, son délai (en log10) et son taux de pertes à partir de la topologie, du routage, des demandes et des poids WFQ réservés aux slices.

L'exploration révèle des pièges qui faussent silencieusement toute modélisation. 42\,\% des chemins du fichier de slices diffèrent du routage simulé, ce qui laisse 8\,\% des flux sans réservation sur leur port d'antenne. Le débit réel vaut 1,25 fois le débit nominal. 93\,\% des flux ont un jumeau aux entrées identiques. Nous les corrigeons ou les documentons avant toute modélisation.

Nous définissons un événement comme un écart brusque au régime habituel du réseau : identificateur de Hampel sur la charge du goulot de chaque flux, environ 14\,\% des flux. Nous comparons, sur des topologies jamais vues, une référence naïve, Ridge, Gradient Boosting, un MLP et un GNN hétérogène flux / files WFQ / liens inspiré de RouteNet.

Le GNN est le plus précis : MAE de 0,175 décade sur le délai (0,277 pour le Gradient Boosting) et de 0,045 sur les pertes (0,067). Surtout, il est le plus robuste pendant les événements : 0,271 contre 0,664. L'ablation montre que ce gain vient de deux sources : les interactions entre flux (sans rétroaction, l'erreur remonte à 0,240) et la représentation explicite des files WFQ par slice (sans elles, 0,428). Nous concluons qu'un GNN apporte une valeur ajoutée quand la QoS dépend de ressources partagées et d'effets en cascade, quand le graphe porte explicitement le levier d'allocation, et surtout quand le réseau sort de son régime habituel.
\end{abstract}
